\documentclass[11pt]{article}

\usepackage[margin=1in]{geometry}
\usepackage{amsmath,amssymb,amsthm,mathtools}
\usepackage{enumitem}
\usepackage{xcolor}
\usepackage{fancyhdr}
\usepackage{hyperref}

\definecolor{navy}{RGB}{24,55,95}
\definecolor{darkgreen}{RGB}{25,100,70}
\hypersetup{colorlinks=true,linkcolor=navy,urlcolor=navy}

\pagestyle{fancy}
\fancyhf{}
\lhead{STAT452 Midterm Review}
\rhead{Regression Proofs}
\cfoot{\thepage}
\setlength{\headheight}{14pt}

\newcommand{\Sxx}{S_{xx}}
\newcommand{\Sxy}{S_{xy}}
\newcommand{\Syy}{S_{yy}}
\newcommand{\bzero}{\widehat{\beta}_0}
\newcommand{\bone}{\widehat{\beta}_1}
\newcommand{\yhat}{\widehat{y}}
\newcommand{\e}{\widehat{e}}
\newcommand{\R}{\mathbb{R}}
\newcommand{\Var}{\operatorname{Var}}
\newcommand{\Cov}{\operatorname{Cov}}
\newcommand{\E}{\operatorname{E}}
\newcommand{\tr}{\operatorname{tr}}

\newenvironment{question}
  {\par\medskip\noindent\color{navy}\bfseries}
  {\par\medskip\color{black}}
\newenvironment{solution}
  {\noindent\color{darkgreen}\textbf{Solution.}\color{black}\ }
  {\hfill$\square$\par\medskip}

\title{Ten Common Proof Questions for Regression\\
       \large Midterm Practice with Complete Solutions}
\author{STAT452 -- Applied Statistics for Engineers and Scientists II}
\date{}

\begin{document}

\maketitle
\tableofcontents
\bigskip

\section*{Setup and notation}
\addcontentsline{toc}{section}{Setup and notation}

Unless a question states otherwise, use the simple linear regression model
\[
  Y_i=\beta_0+\beta_1x_i+\varepsilon_i,\qquad i=1,\ldots,n,
\]
where the predictor values $x_i$ are fixed, $\E(\varepsilon_i)=0$,
$\Var(\varepsilon_i)=\sigma^2$, and the errors are uncorrelated. Define
\[
\bar x=\frac1n\sum_{i=1}^n x_i,\qquad
\bar y=\frac1n\sum_{i=1}^n y_i,
\]
\[
\Sxx=\sum_{i=1}^n(x_i-\bar x)^2,qquad
\Sxy=\sum_{i=1}^n(x_i-\bar x)(y_i-\bar y),
\qquad
\Syy=\sum_{i=1}^n(y_i-\bar y)^2.
\]
We assume $\Sxx>0$, so not all predictor values are equal. The fitted values
and residuals are
\[
  \yhat_i=\bzero+\bone x_i,
  \qquad
  \e_i=y_i-\yhat_i.
\]

\noindent\textbf{Exam strategy.} For most proofs, first write the model and
the relevant estimator in a centered form. Then use one of the key identities
\[
\sum_{i=1}^n(x_i-\bar x)=0,
\qquad
\sum_{i=1}^n(y_i-\bar y)=0,
\qquad
\sum_{i=1}^n \e_i=\sum_{i=1}^n x_i\e_i=0.
\]

\newpage

\section{Derive the least-squares estimators}

\begin{question}
Question 1. Starting from the residual sum of squares
\[
  Q(b_0,b_1)=\sum_{i=1}^n(y_i-b_0-b_1x_i)^2,
\]
prove that its unique minimizers are
\[
  \boxed{\bone=\frac{\Sxy}{\Sxx}},
  \qquad
  \boxed{\bzero=\bar y-\bone\bar x}.
\]
\end{question}

\begin{solution}
Differentiate $Q$ with respect to the two unknown coefficients:
\begin{align*}
\frac{\partial Q}{\partial b_0}
  &=-2\sum_{i=1}^n(y_i-b_0-b_1x_i),\\
\frac{\partial Q}{\partial b_1}
  &=-2\sum_{i=1}^n x_i(y_i-b_0-b_1x_i).
\end{align*}
Setting both derivatives equal to zero gives the normal equations
\begin{align}
  \sum_{i=1}^n y_i&=nb_0+b_1\sum_{i=1}^n x_i, \label{eq:normal1}\\
  \sum_{i=1}^n x_i y_i&=b_0\sum_{i=1}^n x_i+b_1\sum_{i=1}^n x_i^2. \label{eq:normal2}
\end{align}
Divide \eqref{eq:normal1} by $n$:
\[
  \bar y=b_0+b_1\bar x,
\]
so
\[
  b_0=\bar y-b_1\bar x.
\]
Substitute this expression into \eqref{eq:normal2}:
\begin{align*}
\sum x_i y_i
  &=(\bar y-b_1\bar x)\sum x_i+b_1\sum x_i^2\\
  &=n\bar x\bar y+b_1\left(\sum x_i^2-n\bar x^2\right).
\end{align*}
Therefore
\[
 b_1=
 \frac{\sum x_i y_i-n\bar x\bar y}
      {\sum x_i^2-n\bar x^2}.
\]
The numerator and denominator are exactly the centered sums
\[
\sum x_i y_i-n\bar x\bar y=\Sxy,
\qquad
\sum x_i^2-n\bar x^2=\Sxx.
\]
Hence $b_1=\Sxy/\Sxx$ and $b_0=\bar y-b_1\bar x$.

It remains to verify that the stationary point is the unique minimum. The
Hessian matrix is
\[
  2\begin{pmatrix}
      n & \sum x_i\\
      \sum x_i & \sum x_i^2
    \end{pmatrix}.
\]
Its first leading principal minor is $2n>0$, and its determinant is
\[
  4\left(n\sum x_i^2-\left(\sum x_i\right)^2\right)
  =4n\Sxx>0.
\]
Thus the Hessian is positive definite, so the stationary point is the unique
global minimizer.
\end{solution}

\section{Prove equivalent formulas for the slope}

\begin{question}
Question 2. Prove the computational identities
\[
\Sxx=\sum x_i^2-\frac{(\sum x_i)^2}{n},
\qquad
\Sxy=\sum x_iy_i-\frac{(\sum x_i)(\sum y_i)}{n},
\]
and then prove
\[
  \boxed{\bone=r\frac{s_y}{s_x}},
\]
where $r$ is the sample correlation and $s_x,s_y$ are the sample standard
deviations.
\end{question}

\begin{solution}
Expand the definition of $\Sxx$:
\begin{align*}
\Sxx
 &=\sum(x_i-\bar x)^2\\
 &=\sum x_i^2-2\bar x\sum x_i+n\bar x^2\\
 &=\sum x_i^2-2n\bar x^2+n\bar x^2\\
 &=\sum x_i^2-n\bar x^2
 =\sum x_i^2-\frac{(\sum x_i)^2}{n}.
\end{align*}
Similarly,
\begin{align*}
\Sxy
 &=\sum(x_i-\bar x)(y_i-\bar y)\\
 &=\sum x_iy_i-\bar y\sum x_i-\bar x\sum y_i+n\bar x\bar y\\
 &=\sum x_iy_i-n\bar x\bar y\\
 &=\sum x_iy_i-\frac{(\sum x_i)(\sum y_i)}{n}.
\end{align*}
Now use
\[
 r=\frac{\Sxy}{\sqrt{\Sxx\Syy}},
 \qquad
 s_x=\sqrt{\frac{\Sxx}{n-1}},
 \qquad
 s_y=\sqrt{\frac{\Syy}{n-1}}.
\]
Then
\begin{align*}
r\frac{s_y}{s_x}
 &=\frac{\Sxy}{\sqrt{\Sxx\Syy}}
   \sqrt{\frac{\Syy/(n-1)}{\Sxx/(n-1)}}\\
 &=\frac{\Sxy}{\sqrt{\Sxx\Syy}}
   \sqrt{\frac{\Syy}{\Sxx}}
 =\frac{\Sxy}{\Sxx}
 =\bone.
\end{align*}
Thus the least-squares slope equals the correlation multiplied by the ratio
of the sample standard deviations.
\end{solution}

\section{Prove the fundamental residual identities}

\begin{question}
Question 3. For a least-squares line containing an intercept, prove that
\[
 \boxed{\sum_{i=1}^n\e_i=0},
 \qquad
 \boxed{\sum_{i=1}^n x_i\e_i=0},
 \qquad
 \boxed{\sum_{i=1}^n\yhat_i=\sum_{i=1}^n y_i}.
\]
Deduce that $\overline{\widehat y}=\bar y$ and that the fitted line passes
through $(\bar x,\bar y)$.
\end{question}

\begin{solution}
At the least-squares solution, the two normal equations obtained by
differentiating the residual sum of squares are
\[
 \sum_{i=1}^n(y_i-\bzero-\bone x_i)=0
 \quad\text{and}\quad
 \sum_{i=1}^n x_i(y_i-\bzero-\bone x_i)=0.
\]
Because $\e_i=y_i-\bzero-\bone x_i$, these equations are precisely
\[
  \sum_{i=1}^n\e_i=0,
  \qquad
  \sum_{i=1}^n x_i\e_i=0.
\]
Also $y_i=\yhat_i+\e_i$, so summing over $i$ gives
\[
 \sum y_i=\sum\yhat_i+\sum\e_i=\sum\yhat_i.
\]
Dividing by $n$ yields $\overline{\widehat y}=\bar y$. Finally, evaluate the
fitted line at $x=\bar x$:
\[
 \widehat y(\bar x)
 =\bzero+\bone\bar x
 =(\bar y-\bone\bar x)+\bone\bar x
 =\bar y.
\]
Therefore the least-squares line passes through the data centroid
$(\bar x,\bar y)$.
\end{solution}

\section{Prove the regression sum-of-squares decomposition}

\begin{question}
Question 4. Define
\[
 \mathrm{SST}=\sum(y_i-\bar y)^2,
 \quad
 \mathrm{SSR}=\sum(\yhat_i-\bar y)^2,
 \quad
 \mathrm{SSE}=\sum(y_i-\yhat_i)^2.
\]
Prove the ANOVA identity
\[
  \boxed{\mathrm{SST}=\mathrm{SSR}+\mathrm{SSE}}.
\]
State clearly why the cross-product term vanishes.
\end{question}

\begin{solution}
For each observation, add and subtract the fitted value:
\[
 y_i-\bar y=(\yhat_i-\bar y)+(y_i-\yhat_i)
             =(\yhat_i-\bar y)+\e_i.
\]
Square and sum:
\begin{align*}
\sum(y_i-\bar y)^2
 &=\sum(\yhat_i-\bar y)^2+\sum\e_i^2
   +2\sum(\yhat_i-\bar y)\e_i.
\end{align*}
The cross-product term is zero because
\begin{align*}
\sum(\yhat_i-\bar y)\e_i
 &=\sum(\bzero+\bone x_i-\bar y)\e_i\\
 &=(\bzero-\bar y)\sum\e_i+\bone\sum x_i\e_i\\
 &=0,
\end{align*}
where the last line uses both normal-equation identities from Question 3.
Consequently,
\[
 \mathrm{SST}=\mathrm{SSR}+\mathrm{SSE}.
\]
Geometrically, the centered fitted-value vector is orthogonal to the residual
vector, so this is the Pythagorean theorem for the three sums of squares.
\end{solution}

\section{Prove unbiasedness of the coefficient estimators}

\begin{question}
Question 5. Under the fixed-$x$ simple linear regression model, prove that
\[
  \boxed{\E(\bone)=\beta_1},
  \qquad
  \boxed{\E(\bzero)=\beta_0}.
\]
\end{question}

\begin{solution}
Substitute $Y_i=\beta_0+\beta_1x_i+\varepsilon_i$ into the centered slope
formula. Since
\[
 Y_i-\bar Y=\beta_1(x_i-\bar x)+(\varepsilon_i-\bar\varepsilon),
\]
we have
\begin{align*}
\bone
 &=\frac{\sum(x_i-\bar x)(Y_i-\bar Y)}{\Sxx}\\
 &=\frac{\beta_1\sum(x_i-\bar x)^2
   +\sum(x_i-\bar x)(\varepsilon_i-\bar\varepsilon)}{\Sxx}.
\end{align*}
Because $\sum(x_i-\bar x)=0$, the term containing $\bar\varepsilon$
vanishes. Therefore
\[
  \bone=\beta_1+
  \frac{\sum(x_i-\bar x)\varepsilon_i}{\Sxx}.
\]
Taking expectation conditional on the fixed $x_i$ values,
\[
 \E(\bone)
 =\beta_1+\frac{\sum(x_i-\bar x)\E(\varepsilon_i)}{\Sxx}
 =\beta_1.
\]
For the intercept,
\begin{align*}
\bzero
 &=\bar Y-\bone\bar x\\
 &=\beta_0+\beta_1\bar x+\bar\varepsilon-\bone\bar x\\
 &=\beta_0+\bar\varepsilon-\bar x(\bone-\beta_1).
\end{align*}
Thus
\[
 \E(\bzero)
 =\beta_0+\E(\bar\varepsilon)
  -\bar x\{\E(\bone)-\beta_1\}
 =\beta_0.
\]
Normality is not required for either unbiasedness result; zero-mean errors
are sufficient.
\end{solution}

\section{Derive variances, covariance, and fitted-mean variance}

\begin{question}
Question 6. Assuming uncorrelated errors with common variance $\sigma^2$,
prove
\[
 \boxed{\Var(\bone)=\frac{\sigma^2}{\Sxx}},
\]
\[
 \boxed{\Var(\bzero)=\sigma^2\left(\frac1n+
              \frac{\bar x^2}{\Sxx}\right)},
 \qquad
 \boxed{\Cov(\bzero,\bone)=-\frac{\bar x\sigma^2}{\Sxx}}.
\]
Then prove that, at a fixed $x_0$,
\[
 \boxed{\Var\{\widehat{\mu}(x_0)\}
 =\sigma^2\left[\frac1n+\frac{(x_0-\bar x)^2}{\Sxx}\right]},
\]
where $\widehat{\mu}(x_0)=\bzero+\bone x_0$.
\end{question}

\begin{solution}
From Question 5,
\[
 \bone-\beta_1=\sum_{i=1}^n a_i\varepsilon_i,
 \qquad
 a_i=\frac{x_i-\bar x}{\Sxx}.
\]
Using uncorrelated errors with variance $\sigma^2$,
\begin{align*}
\Var(\bone)
 &=\sigma^2\sum a_i^2
 =\sigma^2\frac{\sum(x_i-\bar x)^2}{\Sxx^2}
 =\frac{\sigma^2}{\Sxx}.
\end{align*}
Also
\[
 \bzero-\beta_0=\bar\varepsilon-\bar x(\bone-\beta_1).
\]
We first calculate
\begin{align*}
\Cov(\bar\varepsilon,\bone-\beta_1)
 &=\Cov\left(\frac1n\sum\varepsilon_i,
        \sum a_i\varepsilon_i\right)\\
 &=\frac{\sigma^2}{n}\sum a_i
 =\frac{\sigma^2}{n\Sxx}\sum(x_i-\bar x)
 =0.
\end{align*}
Hence
\begin{align*}
\Var(\bzero)
 &=\Var(\bar\varepsilon)
  +\bar x^2\Var(\bone)
  -2\bar x\Cov(\bar\varepsilon,\bone)\\
 &=\frac{\sigma^2}{n}+\bar x^2\frac{\sigma^2}{\Sxx}.
\end{align*}
Furthermore,
\begin{align*}
\Cov(\bzero,\bone)
 &=\Cov\{\bar\varepsilon-\bar x(\bone-\beta_1),\bone\}\\
 &=0-\bar x\Var(\bone)
 =-\frac{\bar x\sigma^2}{\Sxx}.
\end{align*}
For the fitted mean, it is helpful to rewrite the estimation error in
centered form:
\begin{align*}
\widehat{\mu}(x_0)-(\beta_0+\beta_1x_0)
 &=\bar\varepsilon+(x_0-\bar x)(\bone-\beta_1).
\end{align*}
The two terms on the right are uncorrelated, as shown above. Thus
\begin{align*}
\Var\{\widehat{\mu}(x_0)\}
 &=\Var(\bar\varepsilon)+(x_0-\bar x)^2\Var(\bone)\\
 &=\sigma^2\left[\frac1n+\frac{(x_0-\bar x)^2}{\Sxx}\right].
\end{align*}
For a new individual response $Y_0$, independent of the fitted sample, the
prediction error also contains the new error $\varepsilon_0$. Therefore its
variance is the displayed expression plus $\sigma^2$.
\end{solution}

\section{Prove that MSE is unbiased for the error variance}

\begin{question}
Question 7. In simple linear regression, prove
\[
  \boxed{\E(\mathrm{SSE})=(n-2)\sigma^2}
  \qquad\text{and hence}\qquad
  \boxed{\E(\mathrm{MSE})=\sigma^2},
\]
where $\mathrm{MSE}=\mathrm{SSE}/(n-2)$. You may use matrix notation.
\end{question}

\begin{solution}
Let
\[
 \mathbf Y=\mathbf X\boldsymbol\beta+\boldsymbol\varepsilon,
 \qquad
 \mathbf X=
 \begin{pmatrix}
 1&x_1\\ \vdots&\vdots\\1&x_n
 \end{pmatrix},
\]
where $\mathbf X$ has rank $2$. Define the hat matrix and residual-maker
matrix
\[
 \mathbf H=\mathbf X(\mathbf X^{\mathsf T}\mathbf X)^{-1}
              \mathbf X^{\mathsf T},
 \qquad
 \mathbf M=\mathbf I-\mathbf H.
\]
Because $\mathbf H\mathbf X=\mathbf X$, the residual vector is
\begin{align*}
 \mathbf e
 &=\mathbf Y-\widehat{\mathbf Y}
 =\mathbf M\mathbf Y\\
 &=\mathbf M(\mathbf X\boldsymbol\beta+\boldsymbol\varepsilon)
 =\mathbf M\boldsymbol\varepsilon.
\end{align*}
Both $\mathbf H$ and $\mathbf M$ are symmetric and idempotent. Therefore
\[
 \mathrm{SSE}=\mathbf e^{\mathsf T}\mathbf e
 =\boldsymbol\varepsilon^{\mathsf T}\mathbf M^{\mathsf T}
  \mathbf M\boldsymbol\varepsilon
 =\boldsymbol\varepsilon^{\mathsf T}\mathbf M\boldsymbol\varepsilon.
\]
For any fixed matrix $\mathbf A$ and random vector with mean zero and
covariance $\sigma^2\mathbf I$,
\[
 \E(\boldsymbol\varepsilon^{\mathsf T}\mathbf A
       \boldsymbol\varepsilon)
 =\sigma^2\tr(\mathbf A).
\]
To see this, expand the quadratic form. All off-diagonal expectations are
zero, while each diagonal expectation is $\E(\varepsilon_i^2)=\sigma^2$.
Thus
\[
 \E(\mathrm{SSE})=\sigma^2\tr(\mathbf M)
 =\sigma^2\{\tr(\mathbf I)-\tr(\mathbf H)\}.
\]
Since $\mathbf H$ is the projection onto the two-dimensional column space of
$\mathbf X$, its rank is $2$. An idempotent matrix has eigenvalues only $0$
and $1$, so its trace equals its rank. Consequently,
\[
 \tr(\mathbf M)=n-2
 \quad\Longrightarrow\quad
 \E(\mathrm{SSE})=(n-2)\sigma^2.
\]
Finally,
\[
 \E(\mathrm{MSE})
 =\E\left(\frac{\mathrm{SSE}}{n-2}\right)
 =\sigma^2.
\]
The loss of two degrees of freedom occurs because two coefficients,
$\beta_0$ and $\beta_1$, are estimated from the data.
\end{solution}

\section{Prove the Gauss--Markov result for the slope}

\begin{question}
Question 8. Let $\widetilde\beta_1=\sum_{i=1}^n c_iY_i$ be any linear
unbiased estimator of $\beta_1$. Prove that
\[
  \Var(\widetilde\beta_1)\geq \Var(\bone)
  =\frac{\sigma^2}{\Sxx}.
\]
Thus prove that the least-squares slope is BLUE: the best linear unbiased
estimator of $\beta_1$.
\end{question}

\begin{solution}
First determine the conditions required for unbiasedness. Since
\begin{align*}
\E(\widetilde\beta_1)
 &=\sum c_i(\beta_0+\beta_1x_i)\\
 &=\beta_0\sum c_i+\beta_1\sum c_ix_i,
\end{align*}
this equals $\beta_1$ for every $\beta_0,\beta_1$ exactly when
\[
 \sum c_i=0,
 \qquad
 \sum c_ix_i=1.
\]
The OLS slope is a linear estimator with weights
\[
 a_i=\frac{x_i-\bar x}{\Sxx},
 \qquad
 \bone=\sum a_iY_i.
\]
These weights satisfy
\[
 \sum a_i=0,
 \qquad
 \sum a_ix_i
 =\frac{\sum(x_i-\bar x)x_i}{\Sxx}=1.
\]
For any other linear unbiased estimator, write $c_i=a_i+d_i$. The two
unbiasedness conditions imply
\[
 \sum d_i=0,
 \qquad
 \sum d_ix_i=0.
\]
The OLS weights are orthogonal to the perturbation weights because
\[
 \sum a_id_i
 =\frac1{\Sxx}\sum(x_i-\bar x)d_i
 =\frac1{\Sxx}\left(\sum x_id_i-\bar x\sum d_i\right)
 =0.
\]
Using the common error variance and lack of correlation,
\begin{align*}
\Var(\widetilde\beta_1)
 &=\sigma^2\sum c_i^2\\
 &=\sigma^2\sum(a_i+d_i)^2\\
 &=\sigma^2\sum a_i^2
   +2\sigma^2\sum a_id_i
   +\sigma^2\sum d_i^2\\
 &=\Var(\bone)+\sigma^2\sum d_i^2\\
 &\geq\Var(\bone).
\end{align*}
Equality holds only when all $d_i=0$, meaning the competing estimator is the
OLS estimator itself. This proof does not require normally distributed errors.
\end{solution}

\section{Prove R-squared equals squared correlation and F equals t-squared}

\begin{question}
Question 9. In simple linear regression with an intercept, prove
\[
 \boxed{R^2=r_{xy}^2}.
\]
Also prove that the overall regression $F$ statistic for testing
$H_0:\beta_1=0$ is the square of the usual slope $t$ statistic:
\[
 \boxed{F=t^2}.
\]
\end{question}

\begin{solution}
Because the fitted line passes through $(\bar x,\bar y)$,
\[
 \yhat_i-\bar y
 =\bzero+\bone x_i-\bar y
 =\bone(x_i-\bar x).
\]
Therefore
\begin{align*}
\mathrm{SSR}
 &=\sum(\yhat_i-\bar y)^2
 =\bone^2\sum(x_i-\bar x)^2\\
 &=\left(\frac{\Sxy}{\Sxx}\right)^2\Sxx
 =\frac{\Sxy^2}{\Sxx}.
\end{align*}
Since $\mathrm{SST}=\Syy$,
\[
 R^2=\frac{\mathrm{SSR}}{\mathrm{SST}}
 =\frac{\Sxy^2}{\Sxx\Syy}
 =\left(\frac{\Sxy}{\sqrt{\Sxx\Syy}}\right)^2
 =r_{xy}^2.
\]

For the test of $H_0:\beta_1=0$, the estimated standard error of the slope is
\[
 \operatorname{se}(\bone)=\sqrt{\frac{\mathrm{MSE}}{\Sxx}}.
\]
Thus
\[
 t^2
 =\left(\frac{\bone}{\sqrt{\mathrm{MSE}/\Sxx}}\right)^2
 =\frac{\bone^2\Sxx}{\mathrm{MSE}}
 =\frac{\mathrm{SSR}}{\mathrm{MSE}}.
\]
The regression has one numerator degree of freedom, so its overall statistic is
\[
 F=\frac{\mathrm{MSR}}{\mathrm{MSE}}
 =\frac{\mathrm{SSR}/1}{\mathrm{MSE}}
 =\frac{\mathrm{SSR}}{\mathrm{MSE}}.
\]
Hence $F=t^2$. The tests therefore give identical $p$-values for this
two-sided slope hypothesis.
\end{solution}

\section{Derive the multiple-regression matrix formulas}

\begin{question}
Question 10. Consider the multiple linear regression model
\[
 \mathbf Y=\mathbf X\boldsymbol\beta+\boldsymbol\varepsilon,
 \qquad
 \E(\boldsymbol\varepsilon)=\mathbf 0,
 \qquad
 \Var(\boldsymbol\varepsilon)=\sigma^2\mathbf I_n,
\]
where $\mathbf X$ is an $n\times p$ matrix of full column rank. Prove
\[
 \boxed{\widehat{\boldsymbol\beta}
   =(\mathbf X^{\mathsf T}\mathbf X)^{-1}
     \mathbf X^{\mathsf T}\mathbf Y},
\]
\[
 \boxed{\E(\widehat{\boldsymbol\beta})=\boldsymbol\beta},
 \qquad
 \boxed{\Var(\widehat{\boldsymbol\beta})
   =\sigma^2(\mathbf X^{\mathsf T}\mathbf X)^{-1}}.
\]
Also prove that the hat matrix
\[
 \mathbf H=\mathbf X(\mathbf X^{\mathsf T}\mathbf X)^{-1}
             \mathbf X^{\mathsf T}
\]
is symmetric and idempotent.
\end{question}

\begin{solution}
The residual sum of squares is
\[
 S(\mathbf b)
 =(\mathbf Y-\mathbf X\mathbf b)^{\mathsf T}
   (\mathbf Y-\mathbf X\mathbf b).
\]
Expanding,
\[
 S(\mathbf b)
 =\mathbf Y^{\mathsf T}\mathbf Y
  -2\mathbf b^{\mathsf T}\mathbf X^{\mathsf T}\mathbf Y
  +\mathbf b^{\mathsf T}\mathbf X^{\mathsf T}\mathbf X\mathbf b.
\]
Differentiating with respect to $\mathbf b$ gives
\[
 \frac{\partial S}{\partial\mathbf b}
 =-2\mathbf X^{\mathsf T}\mathbf Y
  +2\mathbf X^{\mathsf T}\mathbf X\mathbf b.
\]
At the minimum,
\[
 \mathbf X^{\mathsf T}\mathbf X\widehat{\boldsymbol\beta}
 =\mathbf X^{\mathsf T}\mathbf Y.
\]
Full column rank makes $\mathbf X^{\mathsf T}\mathbf X$ positive definite
and invertible. Therefore
\[
 \widehat{\boldsymbol\beta}
 =(\mathbf X^{\mathsf T}\mathbf X)^{-1}
   \mathbf X^{\mathsf T}\mathbf Y.
\]
The Hessian is $2\mathbf X^{\mathsf T}\mathbf X$, which is positive
definite, so this stationary point is the unique minimizer.

Substitute the model into the estimator:
\begin{align*}
\widehat{\boldsymbol\beta}
 &=(\mathbf X^{\mathsf T}\mathbf X)^{-1}
   \mathbf X^{\mathsf T}
   (\mathbf X\boldsymbol\beta+\boldsymbol\varepsilon)\\
 &=\boldsymbol\beta
  +(\mathbf X^{\mathsf T}\mathbf X)^{-1}
    \mathbf X^{\mathsf T}\boldsymbol\varepsilon.
\end{align*}
Taking expectations gives
\[
 \E(\widehat{\boldsymbol\beta})=\boldsymbol\beta.
\]
For the covariance matrix, use $\Var(\mathbf A\mathbf Z)
=\mathbf A\Var(\mathbf Z)\mathbf A^{\mathsf T}$:
\begin{align*}
\Var(\widehat{\boldsymbol\beta})
 &=(\mathbf X^{\mathsf T}\mathbf X)^{-1}\mathbf X^{\mathsf T}
   (\sigma^2\mathbf I)
   \mathbf X(\mathbf X^{\mathsf T}\mathbf X)^{-1}\\
 &=\sigma^2(\mathbf X^{\mathsf T}\mathbf X)^{-1}
   \mathbf X^{\mathsf T}\mathbf X
   (\mathbf X^{\mathsf T}\mathbf X)^{-1}\\
 &=\sigma^2(\mathbf X^{\mathsf T}\mathbf X)^{-1}.
\end{align*}
Finally,
\begin{align*}
\mathbf H^{\mathsf T}
 &=\left\{\mathbf X(\mathbf X^{\mathsf T}\mathbf X)^{-1}
             \mathbf X^{\mathsf T}\right\}^{\mathsf T}\\
 &=\mathbf X\left\{(\mathbf X^{\mathsf T}\mathbf X)^{-1}\right\}^{\mathsf T}
   \mathbf X^{\mathsf T}
 =\mathbf H,
\end{align*}
because $\mathbf X^{\mathsf T}\mathbf X$ and its inverse are symmetric. Also
\begin{align*}
\mathbf H^2
 &=\mathbf X(\mathbf X^{\mathsf T}\mathbf X)^{-1}
   \mathbf X^{\mathsf T}\mathbf X
   (\mathbf X^{\mathsf T}\mathbf X)^{-1}\mathbf X^{\mathsf T}\\
 &=\mathbf X(\mathbf X^{\mathsf T}\mathbf X)^{-1}
   \mathbf X^{\mathsf T}
 =\mathbf H.
\end{align*}
Thus $\mathbf H$ is symmetric and idempotent; it is the orthogonal projection
matrix onto the column space of $\mathbf X$.
\end{solution}

\newpage
\section*{One-page proof checklist}
\addcontentsline{toc}{section}{One-page proof checklist}

\begin{enumerate}[leftmargin=*,itemsep=0.55em]
 \item \textbf{OLS coefficients:} Differentiate SSE, obtain the normal
 equations, solve them, and mention positive definiteness for uniqueness.
 \item \textbf{Centering trick:} Use $\sum(x_i-\bar x)=0$ immediately.
 \item \textbf{Residual orthogonality:} With an intercept,
 $\sum\e_i=0$; with a predictor, $\sum x_i\e_i=0$.
 \item \textbf{ANOVA decomposition:} Expand
 $y_i-\bar y=(\yhat_i-\bar y)+\e_i$ and kill the cross term using the normal
 equations.
 \item \textbf{Unbiasedness:} Write the estimator as the true parameter plus
 a weighted sum of zero-mean errors.
 \item \textbf{Variance:} For uncorrelated equal-variance errors,
 $\Var(\sum a_i\varepsilon_i)=\sigma^2\sum a_i^2$.
 \item \textbf{MSE:} Write $\mathrm{SSE}=\boldsymbol\varepsilon^{\mathsf T}
 (\mathbf I-\mathbf H)\boldsymbol\varepsilon$ and use trace $=$ residual
 degrees of freedom.
 \item \textbf{Gauss--Markov:} Express any competing unbiased weights as OLS
 weights plus an orthogonal perturbation.
 \item \textbf{Simple regression identities:}
 $\mathrm{SSR}=\bone^2\Sxx=\Sxy^2/\Sxx$, which quickly gives both
 $R^2=r^2$ and $F=t^2$.
 \item \textbf{Matrix regression:} The normal equations are
 $\mathbf X^{\mathsf T}\mathbf X\widehat{\boldsymbol\beta}
 =\mathbf X^{\mathsf T}\mathbf Y$.
\end{enumerate}

\end{document}
